%=====================================================================
\chapter{The Synopsis and the Research Proposal}\label{ch:synopsis}
%=====================================================================
\locator{6}{Part~6 --- Producing and Defending Research}

\begin{outcomes}
By the end of this chapter you will be able to:
\begin{tight}
  \item \textbf{Assemble} a synopsis in the structure IUB requires, section by
    section.
  \item \textbf{Distinguish} the four sections students most often conflate ---
    problem, rationale, significance and novelty.
  \item \textbf{Navigate} the approval path for your degree, with its
    deadlines.
  \item \textbf{Anticipate} what the Departmental Research Committee and the
    external referee will ask.
\end{tight}
\end{outcomes}

\begin{prereq}
Everything. The synopsis is where Parts I--V are assembled into one document:
the question (\chref{questions}), the review (\chref{slr}), the design
(Part~III), the analysis plan (Part~IV) and the ethics approval
(\chref{ethics}).
\end{prereq}

\begin{keyterms}
synopsis $\bullet$ statement of the problem $\bullet$ rationale $\bullet$
significance $\bullet$ novelty statement $\bullet$ plan of work $\bullet$ Gantt
chart $\bullet$ Departmental Research Committee $\bullet$ external referee
$\bullet$ in-principle approval
\end{keyterms}

\section{What the Synopsis Is For}

It is a contract. It states what you will do, why it is worth doing, and how
you will know whether it worked --- and once approved, it constrains you.
Changing the title afterwards requires going back through the whole approval
chain, which is why the effort spent getting it right is repaid several times
over.

It is also, functionally, a pre-registration (\chref{prereg}) with a committee
attached. Treat it that way and the document gets better.

\begin{regbox}[The approval path]
\textbf{PhD} (\S13). Prepare the synopsis under your supervisor; defend it
before the Departmental Research Committee; on success the DRC recommends it to
AS\&RB through the Board of Studies. It is additionally ``evaluated by at least
one external referee from a national University selected by the Chairperson'',
from ``a list of 6 experts \ldots provided by the Supervisor''. It must be
submitted ``not later than completion of 6\textsuperscript{th} semester \ldots
failing which the topic of the dissertation/thesis will not be approved''.

\medskip
\textbf{MS} (\cite{iub2024mphil}, \S8, \S10). Submit to your supervisor
``within 8 weeks of the start of the 3\textsuperscript{rd} semester'' and defend
before the DRC. The DRC requires coursework completed with CGPA at least
2.70 first. Two copies go to the DRC secretary with a permission letter signed
by the supervisor; a meeting follows within 15 days. The title must be
finalised by the end of the third semester (\S47(iii)).
\end{regbox}

\section{The Structure}

This is the order in the \emph{IUB Theses/Synopses Guidelines-2023}, which is
the format of record for both degrees~\cite{iub2023theses}.

\begin{center}
\small
\begin{longtable}{@{}L{4.0cm}L{10.3cm}@{}}
\toprule
\textbf{Section} & \textbf{What belongs in it} \\
\midrule
\endfirsthead
\toprule
\textbf{Section} & \textbf{What belongs in it} \\
\midrule
\endhead
\bottomrule
\endfoot
Title page
  & Title, scholar, supervisor, department, degree, session. Follow the
    template exactly; deviations are returned \\
Contents, lists, abbreviations
  & Auto-generated. Check they regenerate before every submission \\
Abstract
  & One paragraph: problem, aim, method, expected outcome. Written last \\
Introduction
  & The field, the general problem, and the narrowing to your specific area.
    Ends by pointing at the gap \\
Objectives of the study
  & Numbered, each beginning with a verb, each achievable and checkable.
    Three to five \\
Review of literature
  & Organised by theme, not by paper. Ends with the gap, stated as an absence
    of specific evidence (\chref{slr}) \\
Justification / rationale
  & Why this study, now, here. The argument that it should be done \\
Hypotheses
  & Where the design is confirmatory. Omit, with a sentence saying why, where
    it is exploratory (\chref{questions}) \\
Significance of the study
  & Who benefits, and how, if it succeeds \\
Novelty statement
  & What is new. The delta, in one paragraph (\chref{contribution}) \\
Statement of the problem
  & The problem in one or two precise sentences \\
Methodology and plan of work
  & Design, population and sampling, instruments, procedure, analysis plan,
    ethics, and the timeline \\
Likely benefits / outcomes
  & What will exist at the end: findings, artefact, publications \\
Literature cited
  & APA, complete, and matching the in-text citations exactly \\
\end{longtable}
\end{center}

\begin{alertbox}[Four sections, four different jobs]
These are the ones students merge, and merging them is the commonest reason a
synopsis reads as vague. Each answers a different question.

\begin{center}
\small
\begin{tabular}{@{}L{3.0cm}L{4.0cm}L{6.9cm}@{}}
\toprule
\textbf{Section} & \textbf{Answers} & \textbf{Example} \\
\midrule
Statement of the problem
  & What is wrong?
  & Small software houses in South Punjab detect defects late, and rework
    consumes an estimated fifth of project effort \\
Rationale
  & Why study it now?
  & Checklist-based review is cheap and requires no tooling, but no evidence
    exists for its effect under the time constraints these teams work within \\
Significance
  & Who benefits?
  & Team leads gain evidence for a process decision; the regional software
    industry gains a technique requiring no capital investment \\
Novelty
  & What is new?
  & Prior studies of review checklists have been conducted with untimed
    review; this is the first to vary time-boxing as a factor \\
\bottomrule
\end{tabular}
\end{center}
\end{alertbox}

\section{The Sections That Decide the Outcome}

\subsection{Objectives}

\begin{center}
\small
\begin{tabular}{@{}L{6.4cm}L{7.9cm}@{}}
\toprule
\textbf{Weak} & \textbf{Strong} \\
\midrule
``To study the impact of code review on software quality''
  & ``To measure the difference in defects identified between checklist-guided
    and unguided review in time-boxed conditions'' \\
``To explore machine learning for defect prediction''
  & ``To evaluate whether grouped cross-validation changes the reported
    accuracy of published defect prediction models'' \\
``To develop a framework''
  & ``To design and evaluate a federated aggregation method for clinical
    tabular data with hospital-correlated missingness'' \\
\bottomrule
\end{tabular}
\end{center}

The test: could an examiner tell, at the end, whether you achieved it? ``Study''
and ``explore'' cannot be checked; ``measure'', ``compare'', ``design and
evaluate'' can.

\subsection{The gap}

\begin{pitfall}[``Little research has been done in this area'']
This sentence appears in a large fraction of synopses and is the one an
external referee is most likely to challenge --- often by naming three papers
that did it.

\medskip
A gap is specific: a population not studied, a comparison never made, a
condition never varied, a result never replicated, a method never applied where
its assumptions actually hold. And it is \emph{evidenced}, by the systematic
search of \chref{slr}. A PRISMA flow in the synopsis converts an assertion into
a demonstration, and referees notice.
\end{pitfall}

\subsection{Methodology and plan of work}

This is the section the DRC reads most carefully, because it is where
feasibility is visible.

\begin{center}
\small
\begin{tabular}{@{}L{3.4cm}L{10.9cm}@{}}
\toprule
Design           & Named, and justified against the alternatives you rejected
                   (\chref{questions}) \\
Population and sampling & Who, how many, how recruited, and the a-priori power
                   analysis or saturation argument (\chref{sampling}) \\
Instruments      & Attached as an appendix, with validation status
                   (\chref{measurement}) \\
Procedure        & Step by step, in enough detail to be repeated \\
Analysis plan    & The test or model per objective, stated before any data
                   exists (\chref{inference}) \\
Ethics           & Approval route, consent, data handling (\chref{ethics}) \\
Timeline         & A Gantt chart against the semesters remaining to you \\
Resources        & Equipment, access, cost, and who is providing it \\
\bottomrule
\end{tabular}
\end{center}

\begin{conceptbox}[Make the timeline honest, and it becomes an argument]
A Gantt chart showing data collection finishing the week before submission
tells the committee you have no slack, and they will say so. Build in
contingency, show it, and name the dependency you do not control --- ethics
approval, industrial access, a dataset agreement. A plan that names its risks
is more credible than one that shows a smooth path, and the committee's job is
partly to check you have thought about exactly this.
\end{conceptbox}

\begin{templatebox}[C14 --- Synopsis skeleton]
Fill this before writing prose. If a box cannot be filled in a sentence, that
section is not ready.

\medskip
\begin{tabular}{@{}L{4.0cm}L{9.5cm}@{}}
Working title (12--18 words) & \fillline{9.2cm} \\[7pt]
Problem, one sentence        & \fillline{9.2cm} \\[7pt]
Gap, with evidence           & \fillline{9.2cm} \\[7pt]
Objective 1                  & \fillline{9.2cm} \\[7pt]
Objective 2                  & \fillline{9.2cm} \\[7pt]
Objective 3                  & \fillline{9.2cm} \\[7pt]
RQ / hypothesis per objective & \fillline{9.2cm} \\[7pt]
Design, and one rejected alternative & \fillline{9.2cm} \\[7pt]
Population, N, and its justification & \fillline{9.2cm} \\[7pt]
Instrument(s)                & \fillline{9.2cm} \\[7pt]
Analysis per RQ              & \fillline{9.2cm} \\[7pt]
Novelty, one sentence        & \fillline{9.2cm} \\[7pt]
Significance: who benefits   & \fillline{9.2cm} \\[7pt]
SDG / national agenda link (\S14.1) & \fillline{9.2cm} \\[7pt]
Ethics route                 & \fillline{9.2cm} \\[7pt]
Months to completion, with slack & \fillline{9.2cm} \\[7pt]
Biggest risk, and the fallback & \fillline{9.2cm} \\[7pt]
\end{tabular}
\end{templatebox}

\section{Defending It}

\begin{center}
\small
\begin{tabular}{@{}L{5.0cm}L{9.3cm}@{}}
\toprule
\textbf{What they will ask} & \textbf{What answers it} \\
\midrule
``What exactly is new here?''
  & The novelty statement, delivered in one sentence without hedging \\
``Has this not been done by X?''
  & Your search protocol, and a specific statement of how your study differs
    from X \\
``Why this method and not Y?''
  & The rejected alternative from your design memo (P3) and the constraint that
    decided it \\
``Can you finish this in time?''
  & The Gantt chart, with the slack and the named dependency \\
``What if your main hypothesis is not supported?''
  & That the null result is publishable and why --- which requires that you
    powered for a meaningful effect (\chref{sampling}) \\
``How will you analyse this?''
  & The analysis plan, per objective, already written \\
``How does this serve community needs?''
  & \reg{PhD Regs 2024}{\S14.1} --- have the answer ready; it is a regulation,
    not small talk \\
\bottomrule
\end{tabular}
\end{center}

\begin{alertbox}[The question that most often goes badly]
``What if your hypothesis is not supported?''

\medskip
The wrong answer is that you expect it will be. The right answer states the
smallest effect you powered for, and that failing to find an effect of that
size is itself worth reporting because it would tell practitioners the
technique is not worth adopting under these conditions.

\medskip
A scholar who can answer this has designed a study that cannot fail to produce
knowledge. A scholar who cannot has designed one that only works if the world
cooperates --- and committees are experienced at spotting the difference.
\end{alertbox}

\begin{traditions}{the proposal}
\tradEMP{Hypotheses, power, instruments and analysis plan stated in advance.
The synopsis is a pre-registration with a committee attached.}
\tradDSR{Objectives become the artefact's requirements, and the evaluation
criteria are stated before building (\chref{designscience}).}
\tradQUAL{Objectives are framed as areas of inquiry rather than fixed
questions, with the sampling logic and stopping rule given. Say explicitly that
the design is emergent, so that later adaptation is not read as drift.}
\tradFORM{States the conjecture, the model, and the proof strategy;
feasibility is the argument that the strategy can be carried through.}
\end{traditions}

\begin{lab}[Write the synopsis]
\begin{tightnum}
  \item Complete template C14. Every row, one sentence.
  \item Expand into the IUB section structure. Keep the four distinct sections
    distinct.
  \item Attach the PRISMA flow from \chref{slr} as evidence for the gap.
  \item Attach the instrument and the analysis plan from P4 and P5.
  \item Draw the Gantt chart against your actual remaining semesters, with
    slack shown.
  \item Give it to a colleague and ask them the seven defence questions above.
    Note which you answered badly.
\end{tightnum}

\medskip
This is portfolio piece P7, and it is the deliverable that outlives the course.
\end{lab}

\begin{checkpoint}
\begin{tightnum}
  \item Distinguish statement of the problem, rationale, significance and
    novelty with one sentence each from your own work.
  \item Why is ``to explore X'' a weak objective? Rewrite one of yours.
  \item What is your degree's deadline for synopsis submission, and what
    happens if you miss it?
  \item How would you answer ``what if your hypothesis is not supported?''
\end{tightnum}
\end{checkpoint}

\begin{chaptersummary}
\begin{tight}
  \item The synopsis is a contract and a pre-registration. Changing it later
    costs an approval cycle.
  \item PhD: DRC defence, then BoS/BoF and AS\&RB, plus an external referee
    from a panel of six, by end of semester 6. MS: to the supervisor within 8
    weeks of semester 3, DRC after CGPA 2.70, title fixed by end of semester 3.
  \item Follow the IUB section order exactly, and keep problem, rationale,
    significance and novelty as four distinct answers.
  \item Objectives must be checkable: measure, compare, design and evaluate ---
    not study or explore.
  \item ``Little research exists'' is not a gap. Evidence it with a search
    protocol and a PRISMA flow.
  \item The methodology section is where feasibility shows. Include the power
    analysis, the analysis plan and an honest timeline with named risks.
  \item Be able to say what you will report if the hypothesis fails.
\end{tight}
\end{chaptersummary}

\begin{reviewq}
  \item The synopsis functions as a pre-registration but is not public. What
    would be gained and lost by requiring approved synopses to be publicly
    deposited?
  \item \S13 requires PhD synopsis approval by the sixth semester, and \S15
    requires publication after synopsis approval and before submission.
    Analyse the timing pressure this creates and how a scholar should plan
    around it.
  \item An external referee rejects your synopsis on grounds you believe are
    mistaken. What are your options under the regulations, and what would you
    do?
  \item Committees assess feasibility, but a study that is certainly feasible
    may be unambitious. How should a DRC balance the two, and what would you
    want them to ask?
\end{reviewq}
